%!TEX TS-program = lualatex
%!TEX encoding = UTF-8 Unicode
% Personal CV of Canjia Huang
%
% Based on the Awesome CV LaTeX template:
% https://github.com/posquit0/Awesome-CV
%
% Template author:
% Claud D. Park <posquit0.bj@gmail.com>
% http://www.posquit0.com
%
% Template license:
% CC BY-SA 4.0 (https://creativecommons.org/licenses/by-sa/4.0/)
%


%-------------------------------------------------------------------------------
% CONFIGURATIONS
% Build from the repository root: make (or: latexmk cv.tex)
%-------------------------------------------------------------------------------
% A4 paper size by default, use 'letterpaper' for US letter
\documentclass[11pt, a4paper]{awesome-cv}

% Keep a section title from being left alone at the bottom of a page.
\usepackage{needspace}

%-------------------------------------------------------------------------------
% OPTIONAL: procedural math garden (florilegium)
% Build this variant with `make garden` (or `latexmk cv-florilegium.tex`), which
% defines \withflorilegium before \input-ting this file. Without it the CV is
% unchanged: no package, no page setup change.
% The page keeps its A4 layout and 18.2cm text column. On A4 that leaves no side
% margin for a garden, so the math theme decorates the blank space under the last
% page's text instead (\florilegiumtextend at the end of this file).
%-------------------------------------------------------------------------------
\ifdefined\withflorilegium
  \usepackage[grown, theme=math, density=0.9, intensity=0.9]{florilegium}
  % florilegium 2.0.1 grows in the wider side margin and draws the same garden
  % on every page. This CV keeps its 18.2cm text column, so there is no wide
  % margin: the blank space on page 1 is the two pockets either side of the
  % header, and elsewhere it is the free space under the last page's text.
  % Three small patches, pinned to the vendored 2.0.1:
  %  - geometry(): the garden's own column moves to the paper edge (in a 1.4cm
  %    margin it would hold only stray glyphs), its copper head rule goes
  %    off-page;
  %  - regions(): a page with no free text space gets no footer band, whose few
  %    dots would otherwise land on the footer line;
  %  - field(): add the two header pockets on page 1 (\count0 is the sheet being
  %    shipped). Each is listed twice: the theme sizes its fill from the regions'
  %    area, so twice the area means twice the motifs in the same two pockets.
  % (No Lua -- comments or # operators inside \directlua: TeX would eat them.)
  \directlua{
    local flo = florilegium
    local grow = require("florilegium-grow")
    local lastg
    local geometry = flo.geometry
    flo.geometry = function()
      local g = geometry()
      g.colright = 0.3
      g.headrule = g.ph + 5
      lastg = g
      return g
    end
    local regions = grow.regions
    grow.regions = function(g, layout)
      local rects, col, foot, spill, fx1 = regions(g, layout)
      if spill == nil then
        foot = { foot[1], foot[2], foot[3], foot[2] }
      end
      return rects, col, foot, spill, fx1
    end
    local field = grow.field
    grow.field = function(list)
      local g = lastg
      if g and tex.count[0] == 1 then
        local L, R, H = g.textleft, g.textright, g.ph
        for _ = 1, 2 do
          table.insert(list, { L + 0.05, H - 3.05, L + 3.05, H - 0.45 })
          table.insert(list, { R - 3.05, H - 3.05, R - 0.05, H - 0.45 })
        end
      end
      return field(list)
    end
  }
\fi
% Configure page margins with geometry (unchanged: the garden uses the free
% space below the last page's text, not a margin column)
\geometry{left=1.4cm, top=.8cm, right=1.4cm, bottom=1.8cm, footskip=.5cm}

% Color for highlights (blue theme)
% Awesome Colors: awesome-emerald, awesome-skyblue, awesome-red, awesome-pink, awesome-orange
%                 awesome-nephritis, awesome-concrete, awesome-darknight
\definecolor{awesome}{HTML}{1565C0}
% \colorlet{awesome}{awesome-skyblue}   % the template's own blue, brighter
% \definecolor{awesome}{HTML}{CA63A8}   % or any other colour

% The garden's ink (with \withflorilegium): the CV's own blue instead of the
% package's Delft blue, so garden and headings share one hue.
\ifdefined\withflorilegium
  \definecolor{ink}{HTML}{1565C0}
  % florilegium 2.0.1 colours the math theme's TeX glyphs with
  % text={rgb,255:red,R;green,G;blue,B}, which xcolor reads as a mix of the
  % colours named red, green and blue. This class redefines those three to
  % pastels (awesome-cv.cls, "Basic colors"), so every glyph would come out
  % teal, purple or red. Restore the primaries the mix needs. Nothing else in
  % this CV uses \color{red|green|blue}.
  \definecolor{red}{rgb}{1,0,0}
  \definecolor{green}{rgb}{0,1,0}
  \definecolor{blue}{rgb}{0,0,1}
\fi

% Colors for text
% Uncomment if you would like to specify your own color
% \definecolor{darktext}{HTML}{414141}
% \definecolor{text}{HTML}{333333}
% \definecolor{graytext}{HTML}{5D5D5D}
% \definecolor{lighttext}{HTML}{999999}
% \definecolor{sectiondivider}{HTML}{5D5D5D}

% Set false if you don't want to highlight section with awesome color
\setbool{acvSectionColorHighlight}{true}

% If you would like to change the social information separator from a pipe (|) to something else
\renewcommand{\acvHeaderSocialSep}{\quad\textbar\quad}

% URLs
\newcommand{\BailinDeng}{\href{https://www.bdeng.me}{Bailin Deng}}
\newcommand{\JuanCao}{\href{https://graphics.xmu.edu.cn/~juancao/}{Juan Cao}}
\newcommand{\ZhongguiChen}{\href{https://graphics.xmu.edu.cn/~zgchen/}{Zhonggui Chen}}
\newcommand{\GraphicsXMU}{\href{https://graphics.xmu.edu.cn/}{Graphics@XMU}}
\newcommand{\ACMTOG}{\href{https://dl.acm.org/journal/tog}{ACM Transactions on Graphics}}
\newcommand{\arXiv}{\href{https://arxiv.org}{arXiv}}


%-------------------------------------------------------------------------------
%	PERSONAL INFORMATION
%	Comment any of the lines below if they are not required
%-------------------------------------------------------------------------------
% Available options: circle|rectangle,edge/noedge,left/right
% Add assets/profile.png, then uncomment to show a photo in the header
% \photo{assets/profile.png}
\name{Canjia}{Huang}
\position{Computer Graphics}
\address{Xiamen University, No. 422 Siming South Road, Siming District, Xiamen, Fujian 361005, China}

% \mobile{(+86) 10-9030-1843}
%\whatsapp{(+82) 10-9030-1843}
\email{huangcanjia0214@gmail.com}
%\dateofbirth{January 1st, 1970}
\homepage{https://canjia-huang.github.io/}
\github{Canjia-Huang}
% \linkedin{posquit0}
% \gitlab{gitlab-id}
% \stackoverflow{SO-id}{SO-name}
% \twitter{@twit}
% \x{x-id}
% \skype{skype-id}
% \reddit{reddit-id}
% \medium{medium-id}
% \kaggle{kaggle-id}
% \hackerrank{hackerrank-id}
% \telegram{telegram-username}
% \googlescholar{googlescholar-id}{name-to-display}
%% \firstname and \lastname will be used
% \googlescholar{googlescholar-id}{}
% \extrainfo{extra information}

% \quote{``Be the change that you want to see in the world."}


%-------------------------------------------------------------------------------
\begin{document}

% Print the header with above personal information
% Give optional argument to change alignment(C: center, L: left, R: right)
\makecvheader

% Print the footer with 3 arguments(<left>, <center>, <right>)
% Leave any of these blank if they are not needed
\makecvfooter
  {\today}
  {Canjia Huang~~~·~~~Curriculum Vitae}
  {\thepage}


%-------------------------------------------------------------------------------
%	CV/RESUME CONTENT
%	Each section is imported separately, open each file in turn to modify content
%-------------------------------------------------------------------------------
%-------------------------------------------------------------------------------
%	SECTION TITLE
%-------------------------------------------------------------------------------
\cvsection{Education}


%-------------------------------------------------------------------------------
%	CONTENT
%-------------------------------------------------------------------------------
\begin{cventries}

%---------------------------------------------------------
  \cventry
    {Ph.D. in Computational Mathematics (Combined Master’s–Ph.D. Program)} % Degree
    {School of Mathematical Sciences, Xiamen University} % Institution
    {Xiamen, China} % Location
    {Sep. 2023 - Present}
    {
      \begin{cvitems}
        \item {Advisored by Prof. \JuanCao~and Prof. \ZhongguiChen.}
        \item {Research direction: Computer Graphics, focusing primarily on mesh processing.}
        \item {Member of the Computer Graphics Group \emph{\GraphicsXMU}.}
      \end{cvitems}
    }

  \cventry
    {B.S. in Information and Computing Science} % Degree
    {School of Mathematical Sciences, Xiamen University} % Institution
    {Xiamen, China} % Location
    {Sep. 2018 - Jun. 2022}
    {
      \begin{cvitems} % Description(s) bullet points
        \item {GPA: 3.26/4.}
      \end{cvitems}
    }
%---------------------------------------------------------
\end{cventries}

%-------------------------------------------------------------------------------
%	SECTION TITLE
%-------------------------------------------------------------------------------
\cvsection{Experience}


%-------------------------------------------------------------------------------
%	CONTENT
%-------------------------------------------------------------------------------
\begin{cventries}

%---------------------------------------------------------
  \cventry
    {Teaching Volunteer (Mathematics \& Science/Arts Teacher)} % Job title
    {Yongning County Minning Middle School} % Organization
    {Yinchuan, China} % Location
    {July. 2022 - Aug. 2023} % Date(s)
    {
      \begin{cvitems} % Description(s) of tasks/responsibilities
        \item {Volunteered as a mathematics and science/arts teacher for 7th and 8th-grade students.}
        \item {Instructed specialized courses in robotics and programming, fostering students' computational thinking and hands-on technical skills.}
      \end{cvitems}
    }

  \cventry
    {Trainee / Participant} % Role / Title
    {Youth Marxism Theory Training Program (Graduate Student Supporting Education Corps Special Program), Xiamen University} % Organization
    {Xiamen, China} % Location
    {Oct. 2021 - Jun. 2022} % Date(s)
    {
      \begin{cvitems} % Description(s)
        \item {Completed the training program organized by Xiamen University for the Graduate Student Supporting Education Corps, focusing on political theory, leadership, and comprehensive quality development.}
      \end{cvitems}
    }

%---------------------------------------------------------
\end{cventries}

%-------------------------------------------------------------------------------
% Everything needed to float a figure beside a publication's abstract.
%
% \rightfig[<width as a fraction of \linewidth>]{<image>}
% Puts an image on the right of the paragraph that follows and lets the text
% wrap around it; the text returns to full width below the image.  wrapfig
% cannot be used here (it refuses to work inside cventry's tabular cell), so
% this macro sets a \parshape instead.  \acvfigindent is the left indent used
% for every line but the first (the first line uses \acvfigfirstindent); set
% both to \leftmargin inside a cvitems item.
% The paragraph must be closed with \cvpar
% (see below), otherwise the shape is already forgotten when TeX breaks the
% paragraph into lines.
%-------------------------------------------------------------------------------
% \includegraphics comes from graphicx, which awesome-cv.cls loads via tcolorbox/tikz.
\newbox   \acvfigbox
\newdimen \acvfigwd
\newdimen \acvfigtextwd
\newcount \acvfiglines
\newcount \acvfigcnt
\newdimen \acvfigindent
\setlength{\acvfigindent}{0pt}% left indent of every line but the first
\newdimen \acvfigfirstindent
\setlength{\acvfigfirstindent}{0pt}% left indent of the first line
\newcommand{\acvfigps}{}
% append <n> "<indent> <text width>" pairs to \acvfigps
\def\acvfigbuild#1{%
  \ifnum#1>0\relax
    \edef\acvfigps{\acvfigps\the\acvfigindent\the\acvfigtextwd\space}%
    \acvfigbuild{\number\numexpr#1-1\relax}%
  \fi}
\newcommand{\rightfig}[2][0.30]{%
  \setlength{\acvfigwd}{#1\linewidth}%
  \sbox\acvfigbox{\includegraphics[width=\acvfigwd]{#2}}%
  % Number of text lines covered by the image.
  \acvfiglines=\dimexpr\ht\acvfigbox+\dp\acvfigbox\relax
  \acvfigcnt=\baselineskip
  \acvfiglines=\numexpr(\acvfiglines+\acvfigcnt-1)/\acvfigcnt\relax
  \setlength{\acvfigtextwd}{\dimexpr\linewidth-\acvfigindent-\acvfigwd-1.2em\relax}%
  % Shorten that many lines, then go back to the full width.  The first line
  % keeps its own indent (so a bullet there can start at the left margin).
  \def\acvfigps{\acvfigfirstindent
    \dimexpr\linewidth-\acvfigfirstindent-\acvfigwd-1.2em\relax\space}%
  \acvfigcnt=\numexpr\acvfiglines-1\relax
  \acvfigbuild{\acvfigcnt}%
  \parshape\numexpr\acvfiglines+1\relax \acvfigps
    \acvfigindent\dimexpr\linewidth-\acvfigindent\relax
  % Draw the image at the right edge of the text block, taking no space.
  \smash{\rlap{%
    % the line starts at \acvfigfirstindent, so subtract it to end at the right margin
    \hspace*{\dimexpr\linewidth-\acvfigfirstindent-\acvfigwd\relax}%
    \raisebox{\dimexpr-\ht\acvfigbox\relax}[0pt][0pt]{\usebox\acvfigbox}}}%
}
% \par in disguise: cventry's arguments are not \long, so a plain \par inside
% the description would abort the argument scan.
\newcommand{\cvpar}{\par}


%-------------------------------------------------------------------------------
%	SECTION TITLE
%-------------------------------------------------------------------------------
\cvsection{Publications}

%-------------------------------------------------------------------------------
%	CONTENT
%-------------------------------------------------------------------------------
\begin{cventries}

%---------------------------------------------------------
  \cventry
    {\textbf{Canjia Huang}, \JuanCao\textsuperscript{*}, \BailinDeng, Jiangbo Huang, Ge Yin, \ZhongguiChen}
    {\href{https://dl.acm.org/doi/10.1145/3832781}{Efficient Offset Mesh Generation with Preserved Features and Boundaries}}
    {\ACMTOG}
    {Sep. 14, 2026}
    {
        \begin{cvitems}
        \item \setlength{\acvfigindent}{\leftmargin}%
        \setlength{\acvfigfirstindent}{\leftmargin}%
        \rightfig[0.25]{assets/representation-image-offset-mesh.png}%
        Mesh offsetting plays a pivotal role in CAD/CAM applications. This article introduces an efficient method for generating high-precision, predominantly topologically correct offset meshes from triangular meshes, employing direct offset and point representation. By offsetting sampled points of the input triangular mesh along multiple normals, we create points on a superset of the offset surface. Subsequent feature enhancement captures and enhances features on the point set, facilitating the tailored restricted power diagram-based reconstruction of the offset surface and features. Benchmarking on public datasets, along with comprehensive comparisons with state-of-the-art methods, demonstrates the accuracy, effectiveness, and robustness of our approach. Experimental results showcase its capability to handle diverse models and support arbitrary offset distances, producing predominantly topologically correct offset meshes with small approximation errors while preserving sharp features and open boundaries–even in challenging regions such as thin plates and thin tubular structures.\cvpar
        \end{cvitems}
    }

\cventry
    {Xiaoyang Yu, \textbf{Canjia Huang}, \ZhongguiChen, \JuanCao\textsuperscript{*}}
    {\href{https://arxiv.org/abs/2603.12820}{NeurFrame: Learning Continuous Frame Fields for Structured Mesh Generation}}
    {\arXiv}
    {Mar. 13, 2026}
    {
        \begin{cvitems}
        \item \setlength{\acvfigindent}{\leftmargin}%
        \setlength{\acvfigfirstindent}{\leftmargin}%
        \rightfig[0.25]{assets/representation-image-neurframe.png}%
        Structured meshes, composed of quadrilateral elements in 2D and hexahedral elements in 3D, are widely used in industrial applications and engineering simulations due to their regularity and superior accuracy in finite element analysis. Generating high-quality structured meshes, however, remains challenging, especially for complex geometries and singularities. Field-guided approaches, which construct cross fields in 2D and frame fields in 3D to encode element orientation, are promising but are typically defined on discrete meshes, limiting continuity and computational efficiency. To address these challenges, we introduce \emph{NeurFrame}, a neural framework that represents frame fields continuously over the domain, supporting infinite-resolution evaluation. Trained in a self-supervised manner on discrete mesh samples, NeurFrame produces smooth, high-quality frame fields without relying on dense tetrahedral discretizations. The resulting fields simultaneously guide high-quality quadrilateral surface meshes and hexahedral volumetric meshes, with fewer and better-distributed singularities. By using a single network, NeurFrame also achieves lower computational cost compared to prior self-supervised neural methods that jointly optimize multiple fields.\cvpar
        \end{cvitems}
    }

%---------------------------------------------------------
\end{cventries}

%-------------------------------------------------------------------------------
%	SECTION TITLE
%-------------------------------------------------------------------------------
% Start a new page if the title plus its first entry would not fit,
% so the heading is never left alone at the bottom of a page.
\needspace{7\baselineskip}
\cvsection{Presentation}


%-------------------------------------------------------------------------------
%	CONTENT
%-------------------------------------------------------------------------------
\begin{cventries}

%---------------------------------------------------------
  \cventry
    {Presenter in the Special Session on Intelligent Mesh Technologies} % Role / Title
    {\href{https://graphics.xmu.edu.cn/MEGAS2026/index.html}{MEsh Generation and Applications Symposium (MEGAS 2026)}} % Event
    {Xiamen, China} % Location
    {Nov. 5 - 7, 2026} % Date(s)
    {
      \begin{cvitems} % Description(s)
        \item {Presented the paper \emph{NeurFrame: Learning Continuous Frame Fields for Structured Mesh Generation}.}
      \end{cvitems}
    }

  \cventry
    {Presenter in the Young Scholar Session} % Role / Title
    {\href{https://mp.weixin.qq.com/s/rh_EoJRCUe0eg0RSsuWNpw}{Symposium on Intelligent Grid Methods and Applications (SIGMA 2026)}} % Event
    {Hangzhou, China} % Location
    {May. 29 - Jun. 1, 2026} % Date(s)
    {
      \begin{cvitems} % Description(s)
        \item {Presented the paper \emph{Efficient Offset Mesh Generation with Preserved Features and Boundaries}.}
      \end{cvitems}
    }
%---------------------------------------------------------
\end{cventries}


% With the garden on: let it also fill the free space below the
% text on the last page (must come at the very end of the text).
\ifdefined\withflorilegium
  \florilegiumtextend
\fi

%-------------------------------------------------------------------------------
\end{document}
